\scanpage{97}
\mspage{48}
correspondante de $\mathcal E_{U_{\unsure{\eta|K}}}$ sur $\Gamma'$ --- d'où, en particulier, des actions de $\Gamma'$ sur $\pi$.

(Même chose d'ailleurs pour des sous-corps $k'$ quelconques de $K$, quand on regarde $\mathcal E_K$ comme étant sur $\mathcal E_{k'}$ : la section (29) donne une famille de sections de $\mathcal E_{U_{\unsure{\eta|K}}}$ sur des $\mathcal E_{k'}$, \inserted{ou des sous-groupes ouverts $\Gamma'$ de ceux-ci,} en la composant avec des sections de $\mathcal E_K$ sur des $\mathcal E_{k'}$ --- des sous-groupes $\Gamma'\subset\mathcal E_{k'}$. J'ai envie d'avoir une vue d'ensemble sur cette famille de sections, et de la mettre en relation avec la famille de sections des $\mathcal E_{X_s}$ sur les $\mathcal E_{k(s)}$, \inserted{définie par (29),} définissant des actions \unsure{effectives} (pas seulement ext) des $\mathcal E_{k(s)}$ sur $\pi$. Pour ces dernières, les $k(s)$ n'apparaissent pas comme des sous-corps de $K$, mais plutôt comme des $\underline{\mathcal O}_s/\mathfrak m_s$ --- des corps résiduels de sous-anneaux. Cependant, si on suppose $S$ de dim \deleted{\unsure{finie}} \inserted{\unsure{rel.|relative} 1} sur
\margin{96}
